% GENERATED FILE. Edit canonical lessons, then run npm run generate:books.
% canonical-source-hash: fnv1a64:82ccfbe681c8edea
% canonical-lessons: AR-C309-amir, AR-C309-qasr, AR-C309-taj, AR-C309-arsh, AR-C309-tawwaja

\chapter{Princes and palaces}
\label{ch:ar-princes-and-palaces}
\hlchaptermodality{309}

\begin{chapteropening}
\textbf{By the end of this chapter:} \emph{I can say a prince, a palace, a crown, a throne and to crown.}

Complete the last lesson of chapter 309: I can say a prince, a palace, a crown, a throne and to crown.
\end{chapteropening}

\section[\texorpdfstring{\ar{أمير} (amīr)}{أمير (amīr)}]{\ar{أمير} (amīr) — a prince; an emir}

\label{lesson:AR-C309-amir}



\begin{quote}
Before the new one: say the Arabic for evidence, then the Arabic for a fingerprint.
\end{quote}

\subsection*{You'll want to know: \ar{أمير}}
\textbf{\ar{أمير}} — \emph{amīr} — \textquotedblleft{}a prince; an emir\textquotedblright{}.

\begin{quote}
\ar{كان} \ar{في} \ar{المدينة} \ar{أمير} \ar{شاب}.
\end{quote}

\begin{quote}
\emph{kāna fī al-madīna amīr shābb.}
\end{quote}

\begin{quote}
There was a young prince in the city.
\end{quote}

Plural: \ar{أمراء}, \emph{umarāʾ}.

The feminine is \ar{أميرة}, a princess.

\begin{cousinweb}
\textbf{\ar{أمير}} (\emph{amīr}) is built on the native Arabic root \emph{ʾ-m-r} (Lane), \textquotedblleft{}to command\textquotedblright{}, the root of \emph{amr}, an order. An \emph{amīr} is one who commands; English emir and admiral are usually traced to it.
\end{cousinweb}

\begin{grammarlens}[title={What you've built}]
One more word for telling what happened.
\end{grammarlens}

\subsection*{Your turn}
\begin{itemize}
\raggedright
  \item \emph{Say it:} \emph{amīr}
  \item \emph{Say it:} \emph{amīr}, once more
  \item \emph{Say it:} \emph{kāna fī al-madīna amīr shābb.} — There was a young prince in the city.
  \item \emph{Recall:} say the Arabic for evidence, then the Arabic for a fingerprint, then say \emph{amīr} again
\end{itemize}

\begin{tcolorbox}[breakable,colback=teal!4,colframe=teal!35!black,title={Before you move on}]
What does \ar{أمير} mean? (\textquotedblleft{}A prince; an emir\textquotedblright{}.) Say it once more.
\end{tcolorbox}
\section[\texorpdfstring{\ar{قصر} (qaṣr)}{قصر (qaṣr)}]{\ar{قصر} (qaṣr) — a palace}

\label{lesson:AR-C309-qasr}



\begin{quote}
Before the new one: say the Arabic for a fingerprint, then the Arabic for a prince.
\end{quote}

\subsection*{You'll want to know: \ar{قصر}}
\textbf{\ar{قصر}} — \emph{qaṣr} — \textquotedblleft{}a palace\textquotedblright{}.

\begin{quote}
\ar{زرنا} \ar{قصرا} \ar{قديما} \ar{قرب} \ar{البحر}.
\end{quote}

\begin{quote}
\emph{zurnā qaṣran qadīman qurba al-baḥr.}
\end{quote}

\begin{quote}
We visited an old palace near the sea.
\end{quote}

Plural: \ar{قصور}, \emph{quṣūr}.

\begin{cousinweb}
\textbf{\ar{قصر}} (\emph{qaṣr}) is filed by Lane under the root \emph{q-ṣ-r}, \textquotedblleft{}to confine, to keep within\textquotedblright{}. It was borrowed from Aramaic \emph{qaṣrā} (EtymDB), which is usually traced to Latin \emph{castra}, a camp; the Latin step is not settled.
\end{cousinweb}

\begin{grammarlens}[title={What you've built}]
One more word for telling what happened.
\end{grammarlens}

\subsection*{Your turn}
\begin{itemize}
\raggedright
  \item \emph{Say it:} \emph{qaṣr}
  \item \emph{Say it:} \emph{qaṣr}, once more
  \item \emph{Say it:} \emph{zurnā qaṣran qadīman qurba al-baḥr.} — We visited an old palace near the sea.
  \item \emph{Recall:} say the Arabic for a fingerprint, then the Arabic for a prince, then say \emph{qaṣr} again
\end{itemize}

\begin{tcolorbox}[breakable,colback=teal!4,colframe=teal!35!black,title={Before you move on}]
What does \ar{قصر} mean? (\textquotedblleft{}A palace\textquotedblright{}.) Say it once more.
\end{tcolorbox}
\section[\texorpdfstring{\ar{تاج} (tāj)}{تاج (tāj)}]{\ar{تاج} (tāj) — a crown}

\label{lesson:AR-C309-taj}



\begin{quote}
Before the new one: say the Arabic for a prince, then the Arabic for a palace.
\end{quote}

\subsection*{You'll want to know: \ar{تاج}}
\textbf{\ar{تاج}} — \emph{tāj} — \textquotedblleft{}a crown\textquotedblright{}.

\begin{quote}
\ar{لبست} \ar{الملكة} \ar{تاجا} \ar{من} \ar{الذهب}.
\end{quote}

\begin{quote}
\emph{labisat al-malika tājan min adh-dhahab.}
\end{quote}

\begin{quote}
The queen wore a crown of gold.
\end{quote}

Plural: \ar{تيجان}, \emph{tījān}.

\begin{cousinweb}
\textbf{\ar{تاج}} (\emph{tāj}) was borrowed from Middle Iranian (Parthian) \emph{tāg}, \textquotedblleft{}a crown\textquotedblright{} (EtymDB), an old borrowing; Aramaic \emph{tāgā} is its cousin. Arabic made a verb from it: \emph{tawwaja}, to crown.
\end{cousinweb}

\begin{grammarlens}[title={What you've built}]
One more word for telling what happened.
\end{grammarlens}

\subsection*{Your turn}
\begin{itemize}
\raggedright
  \item \emph{Say it:} \emph{tāj}
  \item \emph{Say it:} \emph{tāj}, once more
  \item \emph{Say it:} \emph{labisat al-malika tājan min adh-dhahab.} — The queen wore a crown of gold.
  \item \emph{Recall:} say the Arabic for a prince, then the Arabic for a palace, then say \emph{tāj} again
\end{itemize}

\begin{tcolorbox}[breakable,colback=teal!4,colframe=teal!35!black,title={Before you move on}]
What does \ar{تاج} mean? (\textquotedblleft{}A crown\textquotedblright{}.) Say it once more.
\end{tcolorbox}
\section[\texorpdfstring{\ar{عرش} (ʿarsh)}{عرش (ʿarsh)}]{\ar{عرش} (ʿarsh) — a throne}

\label{lesson:AR-C309-arsh}



\begin{quote}
Before the new one: say the Arabic for a palace, then the Arabic for a crown.
\end{quote}

\subsection*{You'll want to know: \ar{عرش}}
\textbf{\ar{عرش}} — \emph{ʿarsh} — \textquotedblleft{}a throne\textquotedblright{}.

\begin{quote}
\ar{جلس} \ar{الملك} \ar{على} \ar{العرش}.
\end{quote}

\begin{quote}
\emph{jalasa al-malik ʿalā al-ʿarsh.}
\end{quote}

\begin{quote}
The king sat on the throne.
\end{quote}

Plural: \ar{عروش}, \emph{ʿurūsh}.

\begin{cousinweb}
\textbf{\ar{عرش}} (\emph{ʿarsh}) is built on the native Arabic root \emph{ʿ-r-sh} (Lane), where \emph{ʿarasha} is \textquotedblleft{}to build a roof or a trellis\textquotedblright{}. A raised seat under a canopy became the throne.
\end{cousinweb}

\begin{grammarlens}[title={What you've built}]
One more word for telling what happened.
\end{grammarlens}

\subsection*{Your turn}
\begin{itemize}
\raggedright
  \item \emph{Say it:} \emph{ʿarsh}
  \item \emph{Say it:} \emph{ʿarsh}, once more
  \item \emph{Say it:} \emph{jalasa al-malik ʿalā al-ʿarsh.} — The king sat on the throne.
  \item \emph{Recall:} say the Arabic for a palace, then the Arabic for a crown, then say \emph{ʿarsh} again
\end{itemize}

\begin{tcolorbox}[breakable,colback=teal!4,colframe=teal!35!black,title={Before you move on}]
What does \ar{عرش} mean? (\textquotedblleft{}A throne\textquotedblright{}.) Say it once more.
\end{tcolorbox}
\section[\texorpdfstring{\ar{توج} (tawwaja)}{توج (tawwaja)}]{\ar{توج} (tawwaja) — to crown}

\label{lesson:AR-C309-tawwaja}



\begin{quote}
Before the new one: say the Arabic for a crown, then the Arabic for a throne.
\end{quote}

\subsection*{You'll want to know: \ar{توج}}
\textbf{\ar{توج}} — \emph{tawwaja} — \textquotedblleft{}to crown\textquotedblright{}.

\begin{quote}
\ar{توج} \ar{الأمير} \ar{ملكا} \ar{بعد} \ar{أبيه}.
\end{quote}

\begin{quote}
\emph{tuwwija al-amīr malikan baʿda abīhi.}
\end{quote}

\begin{quote}
The prince was crowned king after his father.
\end{quote}

Present: \ar{يتوج}, \emph{yutawwij}.

In sport: \ar{توج} \ar{الفريق} \ar{بالكأس}, the team won the cup.

\begin{cousinweb}
\textbf{\ar{توج}} (\emph{tawwaja}) comes from \emph{tāj}, \textquotedblleft{}a crown\textquotedblright{}, an Iranian loan (EtymDB). Arabic gave it the Form II pattern: \textquotedblleft{}to put a crown on\textquotedblright{}.
\end{cousinweb}

\begin{grammarlens}[title={What you've built}]
One more word for telling what happened.
\end{grammarlens}

\subsection*{Your turn}
\begin{itemize}
\raggedright
  \item \emph{Say it:} \emph{tawwaja}
  \item \emph{Say it:} \emph{tawwaja}, once more
  \item \emph{Say it:} \emph{tuwwija al-amīr malikan baʿda abīhi.} — The prince was crowned king after his father.
  \item \emph{Recall:} say the Arabic for a crown, then the Arabic for a throne, then say \emph{tawwaja} again
\end{itemize}

\begin{tcolorbox}[breakable,colback=teal!4,colframe=teal!35!black,title={Before you move on}]
What does \ar{توج} mean? (\textquotedblleft{}To crown\textquotedblright{}.) Say it once more.
\end{tcolorbox}
